\chapter*{Résumé}
\addcontentsline{toc}{chapter}{Résumé}

Le présent rapport s'inscrit dans le cadre du module JEE (Java Enterprise Edition). Il décrit la conception et le développement d'une application web de \textbf{gestion des consultations médicales}, dont l'objectif est de centraliser, au sein d'un même outil, la gestion des patients, des médecins et des consultations d'un cabinet médical.

L'application offre deux espaces distincts : un espace \textbf{administrateur} permettant de gérer les dossiers des patients et des médecins, de planifier et de suivre les consultations, et de disposer d'un tableau de bord ; et un espace \textbf{médecin} permettant à chaque praticien de consulter son planning et son historique, d'établir un diagnostic et de rédiger des prescriptions.

Sur le plan technique, l'application repose sur une architecture en couches (présentation, contrôleur, service et accès aux données) et mobilise les technologies JEE : les servlets et les JSP pour la couche web, JPA/Hibernate pour la persistance objet-relationnelle et MySQL pour la base de données. La modélisation du système a été réalisée avec le langage UML, et la gestion du projet s'est appuyée sur une démarche agile de type Scrum, matérialisée par un product backlog et une planification en sprints.

\textbf{Mots-clés} : JEE, JPA, Hibernate, MySQL, Servlets, JSP, UML, architecture en couches, gestion des consultations.

\newpage
